\subsection{Suffix Array}

Seja $s$ uma string de tamanho $n$. O i-ésimo sufixo de $s$ é a substring $s[i \dots n-1]$~. Uma array de sufixos contém inteiros que representam os índices iniciais de todos os sufixos de uma determinada string após os sufixos serem ordenados. Por exemplo, a suffix array de $s = abaab$ é $(2,3,0,4,1)$.

Vamos primeiro fazer um algoritmo para ordenar os shifts cíclicos de uma string. A partir dele, se ordenarmos os shifts cíclicos de uma string $s$ concatenada a um caractere estranho $ \$$, teremos a nossa suffix array em $O(n \log n)$.

\begin{lstlisting}[language=c++, breaklines=true]
vector<int> sort_cyclic_shifts(string const& s) {
    int n = s.size();
    const int alphabet = 256;
    vector<int> p(n), c(n), cnt(max(alphabet, n), 0);
    for (int i = 0; i < n; i++)
        cnt[s[i]]++;
    for (int i = 1; i < alphabet; i++)
        cnt[i] += cnt[i-1];
    for (int i = 0; i < n; i++)
        p[--cnt[s[i]]] = i;
    c[p[0]] = 0;
    int classes = 1;
    for (int i = 1; i < n; i++) {
        if (s[p[i]] != s[p[i-1]])
            classes++;
        c[p[i]] = classes - 1;
    }
    vector<int> pn(n), cn(n);
    for (int h = 0; (1 << h) < n; ++h) {
        for (int i = 0; i < n; i++) {
            pn[i] = p[i] - (1 << h);
            if (pn[i] < 0)
                pn[i] += n;
        }
        fill(cnt.begin(), cnt.begin() + classes, 0);
        for (int i = 0; i < n; i++)
            cnt[c[pn[i]]]++;
        for (int i = 1; i < classes; i++)
            cnt[i] += cnt[i-1];
        for (int i = n-1; i >= 0; i--)
            p[--cnt[c[pn[i]]]] = pn[i];
        cn[p[0]] = 0;
        classes = 1;
        for (int i = 1; i < n; i++) {
            pair<int, int> cur = {c[p[i]], c[(p[i] + (1 << h)) % n]};
            pair<int, int> prev = {c[p[i-1]], c[(p[i-1] + (1 << h)) % n]};
            if (cur != prev)
                ++classes;
            cn[p[i]] = classes - 1;
        }
        c.swap(cn);
    }
    return p;
}
\end{lstlisting}

Isso requer tempo $O(n \log n)$ e $O(n)$ de memória.

A partir dele:

\begin{lstlisting}[language=c++, breaklines=true]
vector<int> suffix_array_construction(string s) {
    s += "$";
    vector<int> sorted_shifts = sort_cyclic_shifts(s);
    sorted_shifts.erase(sorted_shifts.begin());
    return sorted_shifts;
}
\end{lstlisting}

\prob{1}{suffix-array-1} Encontre o menor cyclic shift.

\textbf{Solução:} O algoritmo colocado ordena todos os shifts cíclicos e então $p[0]$ dá a posição do menor shift cíclico.

\prob{2}{suffix-array-2} Ache uma string $s$ dentro de um texto $t$ online. Sabemos o texto $t$ de antemão mas não a string $s$.

\textbf{Solução:} Fazemos a suffix array para o texto $t$ em $O(|t| \log |t|)$. Certamente $s$ deve ser um prefixo de um sufixo de $t$. Então basta fazer uma busca binária por $s$ em $p$. Com isso, demoramos $O(|s| \log |t|)$. Se ocorrer mais vezes, então vão tar próximas, então faz uma segunda busca binária e acha o número de ocorrências dela.

\prob{3}{suffix-array-3} Queremos saber se uma substring de $s$ é menor que outra em $O(1)$

\textbf{Solução:} Depois de construir em $O(|s| \log |s|)$ a suffix array, podemos fazer em $O(1)$:

\begin{lstlisting}[language=c++, breaklines=true]
int compare(int i, int j, int l, int k) {
    pair<int, int> a = {c[k][i], c[k][(i+l-(1 << k))%n]};
    pair<int, int> b = {c[k][j], c[k][(j+l-(1 << k))%n]};
    return a == b ? 0 : a < b ? -1 : 1;
}
\end{lstlisting}

\prob{4}{suffix-array-4} Dada uma string $s$, queremos o longest common prefixo de dois sufixos que iniciam nas posições $i$ e $j$. (aqui fazemos com $O(|s| \log |s|)$ de memória inicial, mas no problema seguinte fazemos sem).

\textbf{Solução:} Temos que nos lembrar das arrays $c[]$ de cada iteração e daí vem a memória adicional.

\begin{lstlisting}[language=c++, breaklines=true]
//log_n e o log(n) na base 2 arredondado pra baixo
int lcp(int i, int j) {
    int ans = 0;
    for (int k = log_n; k >= 0; k--) {
        if (c[k][i % n] == c[k][j % n]) {
            ans += 1 << k;
            i += 1 << k;
            j += 1 << k;
        }
    }
    return ans;
}
\end{lstlisting}

\prob{5}{suffix-array-5} Mesma coisa só que sem memória adicional.

\textbf{Solução:} Construa a array $\text{lcp}[0 \dots n-2]$ onde $\text{lcp}[i]$ é igual ao tamanho do longest common prefix dos sufixos que iniciam em $p[i]$ e $p[i+1]$. Essa array dirá para nós uma resposta para quaisquer dois sufixos adjacentes dessa string. Suponha que seja o pedido de computar o LCP dos sufixos $p[i]$ e $p[j]$. A resposta será $\text{min}(lcp[i], \dots, lcp[j-1])$. Ou seja será uma range minimum query que tem uma variedade de soluções. Para construir a lcp podemos usar o algoritmo de Kasai

\begin{lstlisting}[language=c++, breaklines=true]
vector<int> lcp_construction(string const& s, vector<int> const& p) {
    int n = s.size();
    vector<int> rank(n, 0);
    for (int i = 0; i < n; i++)
        rank[p[i]] = i;

    int k = 0;
    vector<int> lcp(n-1, 0);
    for (int i = 0; i < n; i++) {
        if (rank[i] == n - 1) {
            k = 0;
            continue;
        }
        int j = p[rank[i] + 1];
        while (i + k < n && j + k < n && s[i+k] == s[j+k])
            k++;
        lcp[rank[i]] = k;
        if (k)
            k--;
    }
    return lcp;
}
\end{lstlisting}

\prob{6}{suffix-array-6} Número de substrings distintas

\textbf{Solução:} Pensamos em quais novas substrings começam na posição $p[0]$ depois em $p[1]$ etc.. Tomamos os sufixos ordenadamente e vemos quais prefixos dão novas substrings.  Para cada sufixo atual $p[i]$ teremos novas substrings para todos seus prefixos, exceto para os prefixos que coincidirem com o sufixo $p[i-1]$. Então pegamos todos os prefixos exceto os $\text{lcp}[i-1]$ primeiros. Então a resposta final é:
$$\sum\limits_{i=0}^{n-1} (n - p[i]) - \sum\limits_{i=0}^{n-2} lcp[i] = \dfrac{n^2 + n}{2} - \sum\limits_{i=0}^{n-2} lcp[i]$$

